\section{ĐỀ SỐ \thesection } %LÊ LỢI %Đề 11
\setcounter{bt}{0}

%%Câu 1.
\begin{bt} Tính hợp lý:
	\
	\begin{listEX}[2]	
		\item $131+(-74)+124+(-31)$;
		\item $(-2)^{2} \cdot 38-5 \cdot(-2)^{2}+67 \cdot(-2)^{2}$.
	\end{listEX} 
	\loigiai{a) $150$ \hfill b) $400$.}
\end{bt}


%%Câu 2.
\begin{bt} Tìm số nguyên $x$, biết
	\begin{listEX}[3]
		\item  $30-x=120$;
		\item $-3 \leq x \leq 2$;
		\item $(x-2) \cdot 8+12=(-84)$.
		\end{listEX}
	\loigiai{
	a) $x=-90$ \hfill b) $x\in \{-3;-2;-1;0;1;2\}$ \hfill c) $x=-10$.
	}
\end{bt}



%%Câu 3.
\begin{bt} Bác Dũng có một mảnh vườn hình chữ nhật với chiều dài 20 mét và chiều rộng là 12 mét. Giữa mảnh vườn bác đào một cái ao hình vuông có cạnh 8 mét để nuôi cá, phần còn lại của mảnh vườn bác dùng để trồng rau.
	\immini{
		\begin{enumerate}
			\item Tính chu vi mảnh vườn.
			\item Tính diện tích trồng rau.
			\item Bác Dũng đi mua hạt giống để trồng rau, biết cần gieo 3 gam hạt giống trong 1 mét vuông đất. Hỏi cần bao nhiêu gam hạt giống để gieo hết phần đất trồng rau?
	\end{enumerate}
	}
	{
	\begin{tikzpicture}[scale=1,line width=0.8pt,line join=round,>=stealth,font=\small]
		\def\s{3.2}
		;
		\draw[fill=white] (0,0) rectangle (20/\s,12/\s);
		\draw[pattern=north east lines,pattern color=lightgray] (0,0) rectangle (20/\s,12/\s);
		\draw[fill=white] (6/\s,2/\s) rectangle (14/\s,10/\s); 
		\draw
		(3/\s,6/\s) node{Trồng rau}
		(17/\s,6/\s) node{Trồng rau}
		(10/\s,6/\s) node{Ao cá}
		(10/\s,2/\s) node[below]{8 m}
		(20/\s,6/\s) node[right]{12 m}
		(10/\s,12/\s) node[above]{20 m}
		;
	\end{tikzpicture}
	}

	
	\loigiai{a) 64 m \hfill b) 176 m$^2$ \hfill c) $528$ gam.
	}
\end{bt}


%%Câu 4.
\begin{bt} Số liệu điểm kiểm tra môn Toán giữa học kỳ I của học sinh tổ 2 lớp 6A được ghi lại trong bảng dữ liệu như sau:
	
	\begin{center}
		\begin{tabular}{|c|c|c|c|c|c|}
			\hline
			8 & 10 & 9 & 9 & 7 & 7 \\
			\hline
			9 & 7 & 8 & 7 & 10 & 9 \\
			\hline
		\end{tabular}
	\end{center}
	\
	\begin{enumerate}
		\item  Em hãy cho biết bảng dữ liệu trên đang điều tra về vấn đề gì?
		\item Hãy lập bảng thống kê tương ứng và cho biết tổ 2 lớp 6 A có bao nhiêu thành viên.
	\end{enumerate}
	
	\loigiai{
		\begin{enumerate}
			\item Bảng dữ liệu trên đang điều tra về điểm kiểm tra môn Toán giữa học kỳ I của học sinh tổ 2 lớp 6A. Các số liệu trong bảng đại diện cho điểm số của từng học sinh trong kỳ kiểm tra này.
			\item 
		\begin{tabular}{|c|c|}
			\hline
			\text{Điểm} & \text{Số học sinh} \\
			\hline
			7 & 4 \\
			8 & 2 \\
			9 & 4 \\
			10 & 2 \\
			\hline
		\end{tabular}
	
		
		Tổ 2 lớp 6A có tổng cộng 12 thành viên.
		\end{enumerate}
	}
\end{bt}


%%Câu 5.
\begin{bt} 
Học sinh của một trường THCS khi xếp hàng $20,\,25,\ 30$ em đều thừa 15 em nhưng khi xếp thành hàng 41 em thì vừa đủ hàng. Hỏi trường đó có bao nhiêu học sinh, biết rằng, số học sinh trường đó chưa đến 1000 em.
	\loigiai{ 615 học sinh.
	}
	
\end{bt}


%%Câu 6.
\begin{bt}
Tính nhanh: $\left(4^{234}+4^{235}+4^{236}\right):\left(4^{232}+4^{233}+4^{234}\right)$.
	
	\loigiai{16.}
		




\end{bt}